\subsection{完全整闭整环}

定义 5. 若整环 A 满足下列条件，则称它是完全整闭的：A 的分式域 K 中每个使得所有幂 \(x^{n}\) （ \(n \geqslant 0\) ）都含于 K 的一个有限生成子 A-模的元素 x 属于 A。

注意，诸 \(x^n\) 含于 K 的一个有限生成子 A-模这一假设也可以表述为：存在非零元素 \(d \in \mathbf{A}\) 使得 \(dx^n \in \mathbf{A}\) 对所有 \(n \geqslant 0\) 成立；因为后一条件意味着 \(x^n \in \mathrm{Ad}^{-1}\) ；反之，若 \((b_i)_{1 \leqslant i \leqslant m}\) 是 K 的元素的有限序列，则存在 \(d \in \mathbf{A}\) 使得 \(db_i \in \mathbf{A}\) 对 \(1 \leqslant i \leqslant m\) 成立，从而 \(d\mathbf{M} \subset \mathbf{A}\) 对由诸 \(b_i\) 生成的 K 的子 A-模 M 成立。

显然完全整闭的整环是整闭的；反之，第 1 节命题 1 的推论表明整闭的 Noether 整环是完全整闭的。* 另一方面，高 ≥ 2 的赋值的环（第 VI 章，§ 4，第 4 节）是整闭的但不是完全整闭的。* 若 (A\_t) 是有相同分式域 K 的完全整闭整环的族，则 A = ∩\_t A\_t 是完全整闭的。因为若 x ∈ K 使得对 A 中某个非零 d，\(dx^n\) 对所有 n \textgreater{} 0 属于 A，则假设蕴涵 x ∈ A\_t 对所有 t 成立，从而 x ∈ A。

命题 14. 设 A 是完全整闭的整环。则每个多项式环 \(\mathbf{A}[\mathbf{X}_1,\ldots ,\mathbf{X}_n]\) （相应地，每个形式幂级数环 \(\mathbf{A}[[\mathbf{X}_1,\dots ,\mathbf{X}_n]]\) ）是完全整闭的。

对 \(n\) 作归纳，只需证明 A{[}X{]} （相应地 A{[}{[}X{]}{]}）是完全整闭的。设 P 是 A{[}X{]} （相应地 A{[}{[}X{]}{]}）的分式域的元素，并设存在非零元素 Q ∈ A{[}X{]} （相应地 Q ∈ A{[}{[}X{]}{]}）使得 QP \(^{m}\) ∈ A{[}X{]} （相应地 QP \(^{m}\) ∈ A{[}{[}X{]}{]}）对每个整数 m ≥ 0 成立。若 K 是 A 的分式域，则 A{[}X{]} （相应地 A{[}{[}X{]}{]}）是 K{[}X{]} （相应地 K{[}{[}X{]}{]}）的子环，而 K{[}X{]} （相应地 K{[}{[}X{]}{]}）是主理想整环（《代数》，第 VII 章，§ 1，第 1 节）从而是整闭的（第 4 节，命题 10）且是 Noether 的（《代数》，第 VIII 章，§ 2，第 3 节），因此是完全整闭的；于是我们已经看到 \(\mathbf{P} \in \mathbf{K}[\mathbf{X}]\) （相应地 \(\mathbf{P} \in \mathbf{K}[[\mathbf{X}]]\) ）。设 \(\mathbf{P} = \sum_{k=0}^{\infty} a_k \mathbf{X}^k\) （ \(a_k \in \mathbf{K}\) ）与 \(\mathbf{Q} = \sum_{k=0}^{\infty} b_k \mathbf{X}^k\) （ \(b_k \in \mathbf{A}\) ），我们用反证法论证，设诸 \(a_k\) 并非都属于 \(\mathbf{A}\) ；于是有一个最小的指标 \(i\) 使 \(a_i \notin \mathbf{A}\) ；若我们写

\(P_{1}=\sum_{k=0}a_{k}X^{k}\in A[X]\) ，则由假设立即得出 \(Q(P-P_{1})^{m}\in A[X]\) （相应地 \(Q(P-P_{1})^{m}\in A[[X]]\) ）对所有 \(m\geqslant0\) 也成立。设 j 是使 \(b_{j}\neq0\) 的最小整数；显然在 \(Q(P-P_{1})^{m}\) 中系数 \(\neq0\) 的最低次项是 \(b_{j}a_{i}^{m}X^{j+m^{i}}\) ，因而 \(b_{j}a_{i}^{m}\in A\) 对所有 \(m\geqslant0\) 成立；但由于 A 是完全整闭的，这蕴涵 \(a_{i}\in A\) ，与假设相反。

命题 15. 设 A 是一个带滤链的环，其滤链是穷竭的，且 A 的每个主理想对于由该滤链定义的拓扑是闭的。若相伴分次环 gr(A)（第 III 章，§2，第 3 节）是完全整闭的整环，则 A 是完全整闭的整环。

设 \((\mathbf{A}_n)_{n\in \mathbf{Z}}\) 是 \(\mathbf{A}\) 上定义的滤链；由于 \(\bigcap_{n\in \mathbf{Z}}\mathbf{A}_n\) 是理想 (0) 的闭包（第 III 章，§ 2，第 5 节），假设首先蕴涵滤链 \((\mathbf{A}_n)\) 是分离的，且由于 \(\operatorname{gr}(\mathbf{A})\) 是整环，\(\mathbf{A}\) 也是整环（第 III 章，§ 2，第 3 节，命题 1 的推论）。设 \(x = b / a\) 是分式域 \(\mathbf{K}\) （\(\mathbf{A}\) 的）的元素（ \(a\in \mathbf{A}, b\in \mathbf{A}\) ），对它存在元素 \(d\neq 0\) （\(\mathbf{A}\) 的），使得 \(dx^n\in \mathbf{A}\) 对所有 \(n\geqslant 0\) 成立。我们必须证明 \(b\in \mathbf{A}a\) ，且由假设理想 \(\mathbf{A}a\) 是闭的，只需证明对所有 \(n\in \mathbf{Z}, b\in \mathbf{A}a + \mathbf{A}_n\) 。由于 \(\mathbf{A}\) 的滤链是穷竭的，存在整数 \(q\in \mathbf{Z}\) 使得 \(b\in \mathbf{A}a + \mathbf{A}_q\) 。因此只需证明关系 \(b\in \mathbf{A}a + \mathbf{A}_m\) 蕴涵 \(b\in \mathbf{A}a + \mathbf{A}_{m + 1}\) 。

设 \(b = ay + z\) ，其中 \(y \in \mathbf{A}\) ，\(z \in \mathbf{A}_m\) 。则由假设 \(dx^n \in \mathbf{A}\) 对所有 \(n \geqslant 0\) 成立，由此立即得到 \(d(x - y)^n \in \mathbf{A}\) 对所有 \(n \geqslant 0\) 成立；换句话说，\(dz^n = a^n t_n\) ，其中 \(t_n \in \mathbf{A}\) 对所有 \(n \geqslant 0\) 成立。我们显然可只限于 \(z \neq 0\) 的情形。设 \(v\) 表示 \(\mathbf{A}\) 上的阶函数（第 III 章，§1，第 2 节），我们写 \(v(d) = n_1\) ，\(v(z) = n_2 \geqslant m\) ，\(v(a) = n_3\) ；设 \(d', z', a'\) 分别是 \(d, z, a\) 在 \(\mathbf{A}_{n_1}/\mathbf{A}_{n_1+1}\) 、\(\mathbf{A}_{n_2}/\mathbf{A}_{n_2+1}\) 、\(\mathbf{A}_{n_3}/\mathbf{A}_{n_3+1}\) 中的像。对所有 \(n \geqslant 0\) ，\(v(dz^n) = n_1 + nn_2\) （第 III 章，§2，第 3 节，命题 1），因而 \(\mathrm{gr}(\mathbf{A})\) 中 \(dz^n\) 的典范像是 \(d'z'^n\) ；同样可看出 \(\mathrm{gr}(\mathbf{A})\) 中 \(a^n t_n\) 的典范像形如 \(a'^n t_n'\) ，其中 \(t_n' \in \mathrm{gr}(\mathbf{A})\) ，且由于 \(a' \neq 0\) ，我们推出对所有 \(n \geqslant 0\) ，\(d'(z'/a')^n \in \mathrm{gr}(\mathbf{A})\) 。\(\mathrm{gr}(\mathbf{A})\) 完全整闭这一假设因此蕴涵存在 \(s' \in \mathrm{gr}(\mathbf{A})\) 使得 \(z' = a's'\) ；把 \(s'\) 分解为齐次元素之和，进一步可看出（由于 \(z'\) 与 \(a'\) 是齐次的）\(s'\) 可假设是齐次的，即某个 \(s \in \mathbf{A}\) 的像；于是 \(v(as) = v(z) = n_2\) 且 \(z \equiv as (\mathrm{mod. A_{n_2+1}})\) ；由于 \(n_2 \geqslant m\) ，更有 \(z \equiv as (\mathrm{mod. A_{m+1}})\) ，从而 \(b \equiv a(y + s) (\mathrm{mod. A_{m+1}})\) 。
% semantic-fragments: legacy-input-seam
\relax{}
